\section{Introduction}
\label{sec:intro}

Clear cell renal cell carcinoma (ccRCC) accounts for 70--80\% of renal malignancies and represents the leading cause of kidney cancer mortality, where distant metastasis at presentation is the single most decisive adverse prognostic determinant \cite{schulz2021multimodal, zang2025multimodal, chi2026multimodal}. Conventional prognostic frameworks, including the TNM staging system, Leibovich score, and UISS, rely exclusively on clinicopathological covariates and exhibit modest discriminative capacity for distant metastatic risk \cite{yang2024multimodal}. Because ccRCC metastatic propensity is driven by complex interactions across genomic alterations (\textit{VHL}, \textit{PBRM1}, \textit{SETD2}, \textit{BAP1}) and CT-visible tumor heterogeneity, single-modality predictors fail to fully capture the underlying disease spectrum \cite{schulz2021multimodal, ma2025genomic}.

To capture cross-modal complementarity, multimodal machine learning has emerged as a leading paradigm in biomedical AI \cite{quellec2010dezert, fontani2013dempster, baltrusaitis2018multimodal, huang2020fusion, kline2022multimodal, stahlschmidt2021multimodal}. In oncology, multimodal architectures integrate clinical records, mutation signatures, and radiomic features to improve risk stratification \cite{schulz2021multimodal, mahootiha2023multimodal, yang2024multimodal, zang2025multimodal, chi2026multimodal, xiao2026urofusion}. These reports have driven widespread adoption of information-theoretic decision-level fusion rules, including Bayesian Evidence Fusion, Dempster--Shafer Theory, and Optimal Transport \cite{quellec2010dezert, fontani2013dempster, liu2017weighted, courty2016optimal, shao2023deep, huang2024deep}. Concurrently, retrieval-augmented generation (RAG) has emerged as a powerful paradigm for grounding large language model (LLM) outputs in domain-specific clinical knowledge, yet medical RAG systems that fuse heterogeneous retrieval scores from lexical, dense semantic, and radiomic streams face an analogous score-mixing challenge that has not been formally characterized.

However, both paradigms suffer from a critical methodological limitation: they implicitly assume that base-model probability outputs and retrieval scores are well-calibrated and immune to class-rebalancing distortions. In decision-level fusion, Synthetic Minority Over-sampling Technique (SMOTE) and uncalibrated log-odds stretching distort probability tails, causing information-theoretic fusion rules operating in logit space to manufacture artificial AUROC gains \cite{chen2016calibrating, andersen2024impact, maierhein2024metrics}. In RAG, ad-hoc normalization of heterogeneous retrieval scores---unbounded BM25, bounded dense cosine similarity, and bounded radiomic similarity---produces analogous score-mixing artifacts that inflate retrieval precision and downstream generation faithfulness metrics. Without mandatory out-of-fold probability calibration, reported performance advantages in both paradigms may reflect viewer-side probability distortion rather than genuine clinical complementarity.

We propose \textbf{RenoFusion}\footnote{\url{https://github.com/ali-Hamza817/RenoFusion}}, an open-source two-stage framework that introduces high-performance 3D volumetric multimodal fusion while diagnosing and eliminating the \textbf{Retrieval--Fusion Calibration Hazard} across both decision-level classifier fusion and LLM-powered retrieval-augmented generation for distant metastasis prediction in ccRCC.

\textbf{Stage~I (3D Multimodal Decision-Level Fusion):} Evaluated on an aligned 126-patient TCGA-KIRC cohort ($18~M_1$; 14.29\% prevalence), RenoFusion integrates 3D volumetric CT deep radiomics (88 PyRadiomics + 2,048 MedicalNet 3D ResNet representations), ssGSEA-derived hallmark pathway genomics, and clinical TNM staging. Under strict 5-fold stratified cross-validation with zero data leakage and 2,000 paired bootstrap resamples, 3D volumetric radiomics (M3) elevates single-modality CT discrimination to AUROC 0.8055 [0.7060--0.8982] (surpassing legacy 2D ResNet50 at 0.7037), while clinical TNM staging (M1) achieves AUROC 0.8700 [0.7966--0.9344]. Across seven decision-level fusion architectures, Bayesian Evidential Fusion (BEF) achieves the flagship performance of \textbf{AUROC 0.8756 [0.7991--0.9392]} (AUPRC 0.4989, $F_2$ 0.7143, Brier 0.1052, Platt-calibrated $\text{ECE} = 0.0564$), yielding statistically significant improvements over single-modality imaging ($\Delta\text{AUROC} = +0.0674$, $p = 0.0195$) and outperforming landmark uncalibrated literature benchmarks (Xiao et al., 0.840; Mahootiha et al., 0.800). Crucially, we show that without mandatory Platt cross-calibration, naive rank fusion strategies exhibit severe calibration drift ($\text{ECE} = 0.3611$), whereas calibrated evidential fusion maintains clinical calibration fidelity ($\text{ECE} = 0.0564$).

\textbf{Stage~II (LLM-Powered RAG):} We deploy a real Qwen2.5-3B-Instruct language model with PubMedBERT sentence-transformer dense retrieval across three RAG architectures: (1)~Naive (dense-only), (2)~Multimodal Hybrid (sparse lexical, dense semantic, and radiomic feature fusion), and (3)~Agentic (ReAct tool-calling with structured Thought $\to$ Action $\to$ Observation reasoning traces) on 150 multimodal ccRCC clinical decision queries under 5-fold out-of-fold cross-validation. We show that the Calibration Hazard extends identically to RAG: uncalibrated Reciprocal Rank Fusion inflates retrieval scores ($\text{ECE} = 0.5718$) while calibrated multimodal fusion achieves genuine, statistically significant improvement over naive retrieval (RAGAS composite $+0.0481$, $p = 0.001$), establishing the necessity of natively calibrated retrieval architectures.

The remainder of this paper is organized as follows: Section~\ref{sec:lit_review} provides a chronological literature review and identifies research gaps; Section~\ref{sec:methods} details cohort assembly, base models, fusion rules, RAG architectures, and statistical evaluation; Section~\ref{sec:results} presents experimental findings across both stages; Section~\ref{sec:discussion} discusses mechanistic implications; and Section~\ref{sec:limitations} concludes with study limitations, future directions, and summary conclusions.
